\documentclass[border=8pt,crop=true]{standalone}
\usepackage{circuitikz}
\usetikzlibrary{calc, arrows.meta}
% Shared style for 电子学一 Chapter 2 op-amp circuit figures (CircuiTikZ).
\usetikzlibrary{calc}

\ctikzset{
  bipoles/length=0.95cm,
  bipoles/thickness=1,
  resistors/scale=1,
  capacitors/scale=1,
  label distance=2.5mm,
}

\tikzset{
  opfig/.style={font=\small},
  opfigThin/.style={font=\scriptsize},
}

% Geometry (cm) — keep identical across all op-amp schematics
\def\OpInLen{1.55}    % label → input terminal
\def\OpInGap{0.08}    % terminal → wire to op amp +
\def\OpAmpWire{1.00}  % terminal side → op amp + anchor
\def\OpOutLen{1.05}   % op amp out → output terminal
\def\OpJuncL{0.50}    % v_- → feedback junction (left)
\def\OpFbDown{0.55}   % v_- straight down before R_1
\def\OpRvert{1.20}    % R_1 vertical span to ground
\def\OpInSep{0.55}    % dual-input spacing (v_1 / v_2)
\def\OpFbStub{0.40}   % follower: stub right of output before dropping
\def\OpFbClear{0.22}  % follower: below op-amp south for horizontal feedback


\begin{document}
\begin{circuitikz}[american, scale=1.05, transform shape]
  \coordinate (B) at (0, 0);
  \node[npn, anchor=B, scale=1.05](Q) at (B) {};
  \node[opfig, left=2mm] at (Q.B) {B};
  \node[opfig, below=2mm] at (Q.E) {E};
  \node[opfig, right=2mm] at (Q.C) {C};

  \draw (Q.B) -- ++(-1.35, 0) coordinate (IbIn);
  \draw (Q.E) -- ++(0, -0.95) node[ground, scale=0.85]{};
  \draw (Q.C) -- ++(1.35, 0) coordinate (IcOut);

  \draw[-{Stealth[length=2.2mm]}] ($(IbIn)+(-0.15,0.12)$) -- ++(0.55, 0)
    node[midway, above, opfigThin] {$I_B$};
  \draw[-{Stealth[length=2.2mm]}] ($(Q.C)+(0.15,0.12)$) -- ++(0.55, 0)
    node[midway, above, opfigThin] {$I_C$};
  \draw[-{Stealth[length=2.2mm]}] ($(Q.E)+(0.12,-0.15)$) -- ++(0, -0.55)
    node[midway, right, opfigThin] {$I_E$};

  \node[opfigThin, align=center] at (0, 1.35) {NPN};
\end{circuitikz}
\end{document}
